\section{Summary}
\label{sec:nlp-summary}

The chapter has reported the reference-based NLP metrics at both
document and chapter resolutions
(Sections~\ref{sec:eval-nlp-doc}--\ref{sec:eval-nlp-chapter}) and
the supplementary LLM-as-a-Judge layer (Section~\ref{sec:eval-judge}).
This section brings the document-quality findings together and
frames the questions the Discussion takes up.

The reference-based NLP metrics split the modes into a trade-off
rather than a single ranking. Form leads on recall-oriented
measures, most clearly on BERTScore Recall (Form significantly
higher than both Ask and Chat), with ROUGE-2 and METEOR echoing
it on smaller margins. Chat leads on document-level semantic
alignment, most clearly on SBERT cosine similarity (Chat
significantly higher than both Form and Ask). Form and Chat
trade leadership on BERTScore F1, while Ask underperforms both on
every metric that reaches significance.
The substantive reading is that Form documents recover more of
what the reference \emph{says}, Chat documents stay closer to
what the reference \emph{means}, and Ask documents do neither.

\sloppy
The chapter-level view distributes the document-level trade-off
across chapters rather than across metrics. \emph{Controller and
Contact Details} scores highest across all modes, but this
chapter is filled identically across all three modes by the
user (the LLM has no role in its content); any
apparent mode difference there, including the Ask numerical
lead, is therefore considered user-level noise.
The mode-driven pattern emerges from Chapters 2--6. Form
sweeps the middle chapters (\emph{Purposes and Legal Bases},
\emph{Categories of Personal Data}, \emph{Data Subjects and
Recipients}) on every lexical metric and on BERTScore Recall
and F1. Chat takes the last two chapters (\emph{Retention and
Deletion}, \emph{Technical and Organizational Measures}) on
every metric. The same trade-off that produces Form's
document-level recall edge and Chat's document-level semantic
edge thus appears at the chapter level as a clean two-zone split
across the chapters whose content the modes actually drive.

The supplementary judge layer adds a GDPR-specific
quality reading that the reference\hspace{0pt}-based
metrics cannot directly provide. Form documents are judged the most
complete (Completeness mean $4.84$ of $5$) and the most legally
correct, but the lowest on Scenario Faithfulness ($1.84$) and the
most hallucinated; $83.8\,\%$ of Form documents combine maximum
Completeness with ``moderate'' or worse hallucination, and nearly
half combine maximum Completeness with the floor Hallucination
score. Chat inverts the pattern: most faithful to the study
scenario, least hallucinated, but lowest on Completeness. On the
holistic Overall metric the three modes cluster narrowly ($2.41$
to $2.62$) and all sit below the rubric midpoint; the judge does
not declare a winner.
\fussy

Two tensions surface from the consolidated picture. The mode
participants most often prefer (Form, as reported in
Chapter~\ref{ch:userstudy}) produces the most structurally
complete documents, yet also the most fabricated ones: the
checkbox scaffolding that pushes Completeness toward the ceiling
appears not to constrain what the LLM invents during the final
generation step. And self-report perceived confidence rises from Form to both Ask
and Chat (see Chapter~\ref{ch:userstudy}), with Ask and Chat tied
at the top, while measured document quality follows no such
single ordering: Form leads on recall, Chat leads on
document-level semantic fidelity and on scenario faithfulness,
and the judge declares no clear winner.

With the foundation for RQ1 and RQ2 now set by
Chapters~\ref{ch:userstudy} and~\ref{ch:docquality},
Chapter~\ref{ch:discussion} picks them up and answers them
alongside RQ3 and the further findings that triangulating the
two streams surfaces.